% !TeX root = ../机器学习笔记.tex
% 由主文档通过 \input 载入；请编译主文档。
\setcounter{chapter}{6} % 保留原书章号。
\chapter{线性回归与正则化回归}
\label{chap:ml-regression}

线性回归（linear regression）是从特征预测数值响应的基本模型。它的表达限制清楚、求解结构简单，因此既是实用基线，也是理解估计、正则化与优化的合适起点。

\section{从预测直线到多元线性模型}

模型 $f_{w,b}(x)=w^\top x+b$ 对参数线性；输入可先经过固定变换 $\phi(x)$，于是 $w^\top\phi(x)+b$ 仍是线性参数模型。例如 $\phi(x)=(x,x^2)$ 产生二次曲线。是否使用二次特征属于模型选择，拟合 $w,b$ 属于参数估计。加入特征可以增强表达能力，也会改变样本需求和参数稳定性。

统计分析常写 $y=A\beta+\varepsilon$，其中 $A=[\mathbf1,\mathbf X]$，$\beta=(b,w^\top)^\top$。这个式子只有在同时规定误差条件时才成为具体统计模型。相关性预测不提供因果结论；某个系数的符号反映控制其他所用特征后的拟合关系，并不自动表示改变该特征会改变结果。

\section{最小二乘与参数求解}

训练目标为 $F(\beta)=\frac1{2n}\|A\beta-y\|^2$，梯度是 $A^\top(A\beta-y)/n$。由第\ref{chap:ml-linear-algebra}章的投影结论，正规方程为 $A^\top A\widehat\beta=A^\top y$；满列秩时解唯一，秩亏时可用伪逆得到最小范数解。若先将输入和响应中心化，截距可由 $\widehat b=\bar y-\bar x^\top\widehat w$ 恢复。

对稠密矩阵，计算梯度的成本为 $O(np)$，其中 $p$ 是含截距后的参数数目；$n\geq p$ 时，常用 QR 求解的成本为 $O(np^2)$ 量级。闭式解不是一定更快的实现：当数据大且允许近似解时，迭代方法可避免昂贵分解；当需要高精度且规模适中时，稳定的线性代数求解很有价值。

\section{概率解释与统计性质}

假设 $A$ 固定且满列秩，$\mathbb E[\varepsilon\mid A]=0$、$\operatorname{Cov}(\varepsilon\mid A)=\sigma^2I_n$。由
$\widehat\beta=\beta+(A^\top A)^{-1}A^\top\varepsilon$
直接得到
\[
 \mathbb E[\widehat\beta\mid A]=\beta,\qquad
 \operatorname{Cov}(\widehat\beta\mid A)=\sigma^2(A^\top A)^{-1}.
\]
这两个结论不需要正态性。如果进一步假设误差联合正态，负对数似然关于 $\beta$ 的部分是 $\|y-A\beta\|^2/(2\sigma^2)$，因此极大似然与最小二乘一致，且估计量具有正态分布。

预测一个新特征 $a_*$ 处的均值，方差为 $\sigma^2a_*^\top(A^\top A)^{-1}a_*$；预测具有独立同方差噪声的新响应，还要加上 $\sigma^2$。前者量化均值估计不确定性，后者还包含新观测本身的波动。异方差、相关误差或线性均值假设不成立时，上述方差表达需要修改。

\section{岭回归与收缩}

将数据中心化、不惩罚截距后，岭回归（ridge regression）求解
\begin{equation}
 \min_{w\in\mathbb R^d}\frac1{2n}\|\mathbf Xw-y\|^2
   +\frac\lambda2\|w\|^2,\qquad\lambda>0.
 \label{eq:ml-ridge}
\end{equation}
令梯度为零得到 $\widehat w_\lambda=(\mathbf X^\top\mathbf X+n\lambda I_d)^{-1}\mathbf X^\top y$。因为对任意非零 $v$，
$v^\top(\mathbf X^\top\mathbf X+n\lambda I_d)v=\|\mathbf Xv\|^2+n\lambda\|v\|^2>0$，
解总是唯一，即使特征共线也如此。

设 $\mathbf X=U\Sigma V^\top$，则每个非零奇异方向上的系数由 $\sigma_j/(\sigma_j^2+n\lambda)$ 控制，拟合响应相应乘以 $\sigma_j^2/(\sigma_j^2+n\lambda)$。小奇异值方向收缩更强，减少了对不稳定方向的敏感性，也引入偏差。$\lambda$ 越大并非越好，应在训练内验证选择。

\section{稀疏回归与鲁棒损失}

Lasso（least absolute shrinkage and selection operator）将惩罚换为 $\lambda\|w\|_1$。最优条件是
$0\in\mathbf X^\top(\mathbf Xw-y)/n+\lambda\partial\|w\|_1$；
若一个零系数坐标的损失梯度绝对值不超过 $\lambda$，它可保持为零。弹性网（elastic net）联合使用一、二范数惩罚，兼顾稀疏与收缩。稀疏系数不一定意味着找到了真实因果变量，相关特征可能相互替代。

平方损失对大残差惩罚很强。绝对损失对响应异常值较不敏感，Huber 损失
$\rho_\delta(r)=r^2/2$（$|r|\leq\delta$）和
$\delta|r|-\delta^2/2$（$|r|>\delta$）
则在小残差处光滑拟合，在大残差处线性增长。即使残差损失鲁棒，具有极端输入值的高杠杆点仍可能影响拟合，需分开诊断。

\section{例题、练习与学习检查}

\begin{example}[一维岭收缩]
中心化输入和响应均为 $(-1,0,1)^\top$，取 $\lambda=1/3$。
\end{example}
\begin{solve}
$n=3$，$\mathbf X^\top\mathbf X=\mathbf X^\top y=2$。最小二乘斜率为一，岭斜率为 $2/(2+3\lambda)=2/3$。岭拟合牺牲训练精确性以限制参数，但仅凭这三个点不能断言其总体误差更小。
\end{solve}

习题：推导截距恢复公式；证明岭目标严格凸；令 $\mathbf X^\top\mathbf X/n=I_d$，推导 Lasso 为对 $\mathbf X^\top y/n$ 逐坐标软阈值；比较异常响应和异常输入分别怎样改变最小二乘结果。





